\begin{frame}{One barrier per obligation was not a worse analysis --- it was an unbatched one}
\begin{center}
{\scriptsize
\begin{tabular}{llll}
\toprule
\textbf{formulation} & \textbf{problem shape} & \textbf{result} & \textbf{verdict} \\
\midrule
one fence per obligation & --- & \bad{28--48} barriers; the back end needs 7--10 & the baseline \\
interval greedy & interval cover & splitting a union over-fences & \bad{wrong}, not just larger \\
ILP / exact search & set cover & microseconds at this size & unstable tie-breaks \\
greedy cover & set cover & \good{6 fences for 1284 cross-agent edges} & shipped \\
\bottomrule
\end{tabular}}
\end{center}
\vspace{2mm}
\begin{itemize}
\item A fence is a \key{cover}, not the realization of an obligation: one point discharges every
      edge it happens to stand between. Only cross-agent edges need one, so a single-wave kernel
      gets \good{zero} --- independently agreeing with the back end's own one-wave skip.
\item A loop-carried pair's admissible slots are $\{s>a\}\cup\{s\le b\}$ --- a \key{union wrapping
      the body}, not an interval. That one fact forces set cover.
\item The \emph{points} are minimised; the \emph{names on each point} are not --- token
      permeability, from chapter 6, is what forbids minimising both.
\item ILP was rejected on determinism, not cost --- solver tie-breaking would make the emitted
      assembly non-reproducible.
\end{itemize}
\end{frame}
